\documentclass[12pt]{article}
\usepackage{amsmath, amssymb, amsthm}
\usepackage{geometry}
\usepackage{hyperref}
\usepackage{xcolor}
 
\geometry{margin=1in}
 
% Theorem environments
\newtheorem{theorem}{Theorem}[section]
\newtheorem{lemma}[theorem]{Lemma}
\theoremstyle{definition}
\newtheorem{definition}[theorem]{Definition}
\theoremstyle{remark}
\newtheorem*{remark}{Key Concept}
\newtheorem*{note}{Note}
 
% Custom command for matrix representation
\newcommand{\M}[1]{\mathcal{M}\!\left(#1\right)}
\newcommand{\Mb}[2]{\mathcal{M}\!\left(#1,\,(#2)\right)}
 
\title{\textbf{Axler 3.86: Matrix of Inverse Equals Inverse of Matrix}\\
\large A Proof and Conceptual Commentary}
\author{Linear Algebra Done Right, 4th Edition}
\date{}
 
\begin{document}
 
\maketitle
 
\section*{Context in the Text}
 
Theorem 3.86 is the \textbf{last result stated in Chapter 3D} of \emph{Linear Algebra Done Right}, 4th edition, appearing just before
the end-of-chapter exercises. It is presented without proof, with Axler noting
only that ``the proof of the next result is left as an exercise.''
 
Chapter 3D culminates in the change-of-basis formula (3.84), which shows how
the matrix of an operator transforms under a change of basis via
$A = C^{-1}BC$. Theorem 3.86 is a natural coda to this: having established how
matrices of operators compose and change with basis, the chapter closes by
confirming that the matrix correspondence also respects inversion. It is
intentionally elementary --- a short exercise designed to make the reader
explicitly invoke the multiplicativity machinery developed in the chapter rather
than taking the result for granted.
 
\section*{Statement}
 
\begin{theorem}[3.86 --- Matrix of inverse equals inverse of matrix]
Suppose $v_1, \ldots, v_n$ is a basis of $V$ and $T \in \mathcal{L}(V)$ is
invertible. Then
\[
    \M{T^{-1}} = \bigl(\M{T}\bigr)^{-1},
\]
where both matrices are with respect to the basis $v_1, \ldots, v_n$.
\end{theorem}
 
\section*{Prerequisites}
 
This proof relies on one named result and one direct consequence of the definition:
 
\begin{itemize}
    \item \textbf{Multiplicativity of $\mathcal{M}$} (Axler 3.43): For operators
          $S, T \in \mathcal{L}(V)$ with respect to the same basis,
          \[
              \M{ST} = \M{S}\,\M{T}.
          \]
 
    \item \textbf{Matrix of the identity} (from the definition of $\mathcal{M}$):
          The $k$-th column of $\mathcal{M}(T)$ is determined by writing $Tv_k$
          as a linear combination of $v_1, \ldots, v_n$. For the identity map,
          $Iv_k = v_k$, so the $k$-th column has a $1$ in position $k$ and $0$s
          elsewhere. Therefore $\mathcal{M}(I) = I$ (the $n \times n$ identity
          matrix) follows directly from the definition, without a separate theorem.
\end{itemize}
 
\section*{Proof}
 
Let $B = \M{T}$ and $B' = \M{T^{-1}}$, both with respect to $v_1, \ldots, v_n$.
 
Since $T$ is invertible, we have $TT^{-1} = I$ as operators. Applying
$\mathcal{M}$ to both sides and invoking multiplicativity:
\[
    BB' \;=\; \M{T}\,\M{T^{-1}} \;=\; \M{TT^{-1}} \;=\; \M{I} \;=\; I.
\]
By an identical argument using $T^{-1}T = I$:
\[
    B'B \;=\; \M{T^{-1}}\,\M{T} \;=\; \M{T^{-1}T} \;=\; \M{I} \;=\; I.
\]
Since $BB' = I$ and $B'B = I$, we conclude that $B' = B^{-1}$, i.e.,
\[
    \M{T^{-1}} = \bigl(\M{T}\bigr)^{-1}. \qed
\]
 
\section*{Key Concept: Why This Exercise Exists}
 
\begin{remark}
The result \emph{feels} obvious, and it is --- but only once you have
identified the two bridges being crossed:
 
\begin{enumerate}
    \item \textbf{Operator composition $\longrightarrow$ matrix multiplication.}
          The map $T \mapsto \mathcal{M}(T)$ converts operator composition into
          matrix multiplication. This is non-trivial and was the whole motivation
          behind defining matrix multiplication the way it is in Chapter 3.
 
    \item \textbf{Identity operator $\longrightarrow$ identity matrix.}
          The map sends $I$ (the identity \emph{operator}) to $I$ (the identity
          \emph{matrix}). Without explicitly invoking $\M{I} = I$, one has not
          actually closed the argument.
\end{enumerate}
 
The proof chains these two facts together: the operator equation $TT^{-1} = I$
is lifted to the matrix equation $\M{T}\M{T^{-1}} = I$ by crossing both bridges
simultaneously.
\end{remark}
 
\begin{note}
Algebraically, the map $T \mapsto \mathcal{M}(T)$ (with respect to a fixed
basis) is a \emph{ring isomorphism} from $\mathcal{L}(V)$ to $M_{n \times n}(F)$.
Ring isomorphisms necessarily:
\begin{itemize}
    \item send the multiplicative identity to the multiplicative identity, and
    \item send invertible elements to invertible elements, preserving inverses.
\end{itemize}
Theorem 3.86 is precisely the second of these properties. The ``obviousness''
of the result is a reflection of how well-designed the correspondence between
operators and matrices is --- but the formal proof still requires citing the
machinery that establishes that correspondence.
\end{note}
 
\end{document}